% GENERATED FILE. Edit canonical lessons, then run npm run generate:books.
% canonical-source-hash: fnv1a64:5feaa06771d199a6
% canonical-lessons: HI-C75-first, HI-C75-second, HI-C75-third, HI-C75-fourth, HI-C75-fifth

\chapter{First, Second, Third}
\label{ch:hi-ordinals}
\hlchaptermodality{82}

\begin{chapteropening}
\textbf{By the end of this chapter:} \emph{I can say first to fifth in Hindi with the right ending for the noun, and say why the first four have to be learned as words and the fifth does not.}

The fifth ordinal states the -van rule against the four irregular words that came before it.
\end{chapteropening}

\section[pahlā / pahlī / pahle]{\hi{पहला} (pahlā) — first}

\label{lesson:HI-C75-first}



\begin{quote}
A new kind of counting starts here, so the recalls come first, at the four growing distances this book measures.

\begin{itemize}
\raggedright
  \item \emph{Recall:} say \emph{der} — \textbf{R1}, one lesson back, before it decays
  \item \emph{Recall:} say \emph{khulā} — \textbf{R2}, the first real retrieval, five to fifteen lessons back
  \item \emph{Recall:} say \emph{the two form names, Mīrā and Aruṇ} — \textbf{R3}, durable, twenty to sixty lessons back
  \item \emph{Recall:} say \emph{deś} — \textbf{R3}, a second word from the same distance
  \item \emph{Recall:} say \emph{ghaṇṭā} — \textbf{R4}, recognition at distance, eighty lessons back or more
  \item \emph{Recall:} say \emph{harā} — \textbf{R4}, a second word from the same distance
\end{itemize}
\end{quote}

\subsection*{You'll want to know: \hi{पहला}}
\textbf{\hi{पहला}} (\emph{pahlā}) — \textquotedblleft{}\textbf{first}\textquotedblright{}.

\textbf{\hi{एक}, \hi{दो}, \hi{तीन}} answer \textbf{how many}. This answers \textbf{which one in the line}, and Hindi keeps a second set of words for that.

It agrees like any \emph{-ā} adjective: \textbf{\hi{पहला} \hi{दिन}} (\emph{pahlā din}), the first day, but \textbf{\hi{पहली} \hi{किताब}} (\emph{pahlī kitāb}), the first book — feminine noun, feminine ending.

\begin{cousinweb}
There is no \textbf{\hi{एक}} inside \textbf{\hi{पहला}} (\emph{pahlā}). It comes down from Sanskrit \textbf{\hi{प्रथम}} (\emph{prathama}), a word that already meant \textquotedblleft{}foremost\textquotedblright{} before Hindi existed, and the number was never part of it.

You will meet its close relative constantly: \textbf{\hi{पहले}} (\emph{pahle}) — \textquotedblleft{}\textbf{before}, \textbf{earlier}, \textbf{first of all}\textquotedblright{}. \emph{\hi{खाने} \hi{से} \hi{पहले}} is \textquotedblleft{}before eating\textquotedblright{}. So in Hindi the word for \emph{first} and the word for \emph{before} are the same word wearing a different ending, which is a reasonable thing for a language to decide: putting one thing before another \textbf{is} putting it first.
\end{cousinweb}

\subsection*{Your turn}
Say these aloud:

\begin{itemize}
\raggedright
  \item \emph{pahlā}
  \item \emph{pahlā din} — the first day
  \item \emph{pahlī kitāb} — the first book
\end{itemize}

\begin{tcolorbox}[breakable,colback=teal!4,colframe=teal!35!black,title={Before you move on}]
Say \textquotedblleft{}the first day\textquotedblright{}. (\textbf{\hi{पहला}} (\emph{pahlā}) \hi{दिन}.) What changes in \emph{\hi{पहली} \hi{किताब}}, and why? (\textbf{The ending} — \emph{kitāb} is feminine.) And what does \textbf{\hi{पहले}} mean? (\textbf{Before, earlier} — the same word.)
\end{tcolorbox}
\section[dūsrā / dūsrī]{\hi{दूसरा} (dūsrā) — second, and the other one}

\label{lesson:HI-C75-second}



\begin{quote}
Recalls before the second ordinal, at the four measured distances.

\begin{itemize}
\raggedright
  \item \emph{Recall:} say \emph{pahlā} — \textbf{R1}, one lesson back, before it decays
  \item \emph{Recall:} say \emph{band} — \textbf{R2}, the first real retrieval, five to fifteen lessons back
  \item \emph{Recall:} say \emph{the supported name entry} — \textbf{R3}, durable, twenty to sixty lessons back
  \item \emph{Recall:} say \emph{Bhārat} — \textbf{R3}, a second word from the same distance
  \item \emph{Recall:} say \emph{the letter kha} — \textbf{R4}, recognition at distance, eighty lessons back or more
  \item \emph{Recall:} say \emph{pīlā} — \textbf{R4}, a second word from the same distance
\end{itemize}
\end{quote}

\subsection*{You'll want to know: \hi{दूसरा}}
\textbf{\hi{दूसरा}} (\emph{dūsrā}) — \textquotedblleft{}\textbf{second}\textquotedblright{}.

\textbf{\hi{दूसरा} \hi{दिन}} (\emph{dūsrā din}) — the second day. \textbf{\hi{दूसरी} \hi{किताब}} (\emph{dūsrī kitāb}) — the second book.

Like \emph{pahlā}, there is no \textbf{\hi{दो}} visible in it, though this one does descend from the Sanskrit word for two.

\subsection*{You'll want to know: the other meaning}
\textbf{\hi{दूसरा}} (\emph{dūsrā}) also means \textquotedblleft{}\textbf{other}, \textbf{another}\textquotedblright{}, and that is not a separate word to learn — it is the same one.

\begin{itemize}
\raggedright
  \item \textbf{\hi{दूसरा} \hi{आदमी}} (\emph{dūsrā ādmī}) — the other man.
  \item \textbf{\hi{दूसरा} \hi{घर}} (\emph{dūsrā ghar}) — the other house.
\end{itemize}

Hold both senses together. A chapter two ahead turns on them: once you can say \emph{ek} and \emph{dūsrā} in one breath, you can say \textquotedblleft{}\textbf{one … the other}\textquotedblright{}, which is a thing Hindi does with no extra machinery at all.

\subsection*{Your turn}
Say these aloud:

\begin{itemize}
\raggedright
  \item \emph{dūsrā}
  \item \emph{dūsrā din} — the second day
  \item \emph{dūsrā ādmī} — the other man
\end{itemize}

\begin{tcolorbox}[breakable,colback=teal!4,colframe=teal!35!black,title={Before you move on}]
Say \textquotedblleft{}the second book\textquotedblright{}. (\textbf{\hi{दूसरी} \hi{किताब}}.) What is the other thing \textbf{\hi{दूसरा}} (\emph{dūsrā}) means? (\textbf{Other, another}.) And say \textquotedblleft{}the other man\textquotedblright{}. (\textbf{\hi{दूसरा}} (\emph{dūsrā}) \textbf{\hi{आदमी}}.)
\end{tcolorbox}
\section[tīsrā / tīsrī]{\hi{तीसरा} (tīsrā) — third}

\label{lesson:HI-C75-third}



\begin{quote}
Recalls first, then the third.

\begin{itemize}
\raggedright
  \item \emph{Recall:} say \emph{dūsrā} — \textbf{R1}, one lesson back, before it decays
  \item \emph{Recall:} say \emph{praveś} — \textbf{R2}, the first real retrieval, five to fifteen lessons back
  \item \emph{Recall:} say \emph{the delayed name entry} — \textbf{R3}, durable, twenty to sixty lessons back
  \item \emph{Recall:} say \emph{śahar} — \textbf{R3}, a second word from the same distance
  \item \emph{Recall:} say \emph{umr} — \textbf{R4}, recognition at distance, eighty lessons back or more
  \item \emph{Recall:} say \emph{din} — \textbf{R4}, a second word from the same distance
\end{itemize}
\end{quote}

\subsection*{You'll want to know: \hi{तीसरा}}
\textbf{\hi{तीसरा}} (\emph{tīsrā}) — \textquotedblleft{}\textbf{third}\textquotedblright{}.

\textbf{\hi{तीसरा} \hi{दिन}} (\emph{tīsrā din}), \textbf{\hi{तीसरी} \hi{किताब}} (\emph{tīsrī kitāb}).

Say \textbf{\hi{तीन}} and then \textbf{\hi{तीसरा}} and you can hear the number this time — the first ordinal where you can. And say \emph{dūsrā, tīsrā} together: they rhyme, because they share an ending, \textbf{-\hi{सरा}}, that they share with nothing else in the language.

Two words with a private ending is not a pattern; it is a pair. The next lesson goes a third way, and the one after that finally states a rule.

\subsection*{Your turn}
Say these aloud:

\begin{itemize}
\raggedright
  \item \emph{tīsrā}
  \item the join — \emph{tīn} $\to$ \emph{tīsrā}
  \item the pair — \emph{dūsrā, tīsrā}
\end{itemize}

\begin{tcolorbox}[breakable,colback=teal!4,colframe=teal!35!black,title={Before you move on}]
Say the third day. (\textbf{\hi{तीसरा} \hi{दिन}}.) Which number can you hear inside it? (\textbf{\hi{तीन}}.) And how many words share the \emph{-\hi{सरा}} ending? (\textbf{Two} — \emph{dūsrā} and \emph{tīsrā}.)
\end{tcolorbox}
\section[cauthā / cauthī]{\hi{चौथा} (cauthā) — fourth}

\label{lesson:HI-C75-fourth}



\begin{quote}
Recalls first, then the fourth.

\begin{itemize}
\raggedright
  \item \emph{Recall:} say \emph{tīsrā} — \textbf{R1}, one lesson back, before it decays
  \item \emph{Recall:} say \emph{nikās} — \textbf{R2}, the first real retrieval, five to fifteen lessons back
  \item \emph{Recall:} say \emph{cāhnā} — \textbf{R3}, durable, twenty to sixty lessons back
  \item \emph{Recall:} say \emph{bhāṣā} — \textbf{R3}, a second word from the same distance
  \item \emph{Recall:} say \emph{kitne sāl ke ho?} — \textbf{R4}, recognition at distance, eighty lessons back or more
  \item \emph{Recall:} say \emph{rāt} — \textbf{R4}, a second word from the same distance
\end{itemize}
\end{quote}

\subsection*{You'll want to know: \hi{चौथा}}
\textbf{\hi{चौथा}} (\emph{cauthā}) — \textquotedblleft{}\textbf{fourth}\textquotedblright{}.

\textbf{\hi{चौथा} \hi{दिन}} (\emph{cauthā din}), \textbf{\hi{चौथी} \hi{किताब}} (\emph{cauthī kitāb}).

\textbf{\hi{चार}} is four, and \emph{cauthā} comes from Sanskrit \textbf{\hi{चतुर्थ}} (\emph{caturtha}) — you can hear the four, but the ending is \textbf{-\hi{था}}, not the \textbf{-\hi{सरा}} of the last two.

Count what has happened. Four ordinals, and four different shapes: \emph{pahlā} from a word that never meant one, \emph{dūsrā} and \emph{tīsrā} sharing a private ending, and now \emph{cauthā} with an ending of its own. Nothing so far has generalised.

That is the honest state of it, and it is also the good news: these four are the whole of the memorising. The next lesson states the rule, and it holds for everything after — with one hole, which gets a lesson of its own.

\subsection*{Your turn}
Say these aloud:

\begin{itemize}
\raggedright
  \item \emph{cauthā}
  \item the join — \emph{cār} $\to$ \emph{cauthā}
  \item the four so far — \emph{pahlā, dūsrā, tīsrā, cauthā}
\end{itemize}

\begin{tcolorbox}[breakable,colback=teal!4,colframe=teal!35!black,title={Before you move on}]
Say the fourth day. (\textbf{\hi{चौथा}} (\emph{cauthā}) \textbf{\hi{दिन}}.) What ending does it take? (\textbf{-\hi{था}}, not the \emph{-\hi{सरा}} of the last two.) And how many of the ordinals have to be learned as separate words? (\textbf{Four} — and you now have them.)
\end{tcolorbox}
\section[pā̃cvā̃]{\hi{पाँचवाँ} (pā̃cvā̃) — fifth, and the rule at last}

\label{lesson:HI-C75-fifth}



\begin{quote}
Recalls first, then the fifth — and the ending that ends the memorizing.

\begin{itemize}
\raggedright
  \item \emph{Recall:} say \emph{cauthā} — \textbf{R1}, one lesson back, before it decays
  \item \emph{Recall:} say \emph{der} — \textbf{R2}, the first real retrieval, five to fifteen lessons back
  \item \emph{Recall:} say \emph{the infinitive as an object} — \textbf{R3}, durable, twenty to sixty lessons back
  \item \emph{Recall:} say \emph{aṅgrezī} — \textbf{R3}, a second word from the same distance
  \item \emph{Recall:} say \emph{the letter da} — \textbf{R4}, recognition at distance, eighty lessons back or more
  \item \emph{Recall:} say \emph{śubh rātri} — \textbf{R4}, a second word from the same distance
\end{itemize}
\end{quote}

\subsection*{You'll want to know: \hi{पाँचवाँ}}
\textbf{\hi{पाँचवाँ}} (\emph{pā̃cvā̃}) — \textquotedblleft{}\textbf{fifth}\textquotedblright{}.

\textbf{\hi{पाँचवाँ} \hi{दिन}} (\emph{pā̃cvā̃ din}), \textbf{\hi{पाँचवीं} \hi{किताब}} (\emph{pā̃cvī̃ kitāb}).

Two nasal hooks in one word: one over the \textbf{\hi{पाँ}}, one over the \textbf{\hi{वाँ}}. The chandrabindu you learned in the writing chapters is doing both.

\begin{grammarlens}[title={Grammar Lens: the number, plus -\hi{वाँ}}]
\textbf{\hi{पाँच}} $\to$ \textbf{\hi{पाँचवाँ}}. Nothing is taken away and nothing is reshaped: the cardinal stands whole and \textbf{-\hi{वाँ}} goes on the end.

It goes on doing that for every number above five. \textbf{\hi{बीस}} gives \textbf{\hi{बीसवाँ}}, twentieth; \textbf{\hi{सौ}} gives \textbf{\hi{सौवाँ}}, hundredth. The four words before this one were the exceptions, and you met them first because a learner needs \emph{first} and \emph{second} on the first day and cannot wait for a rule.

One number after this breaks it. Everything else obeys.
\end{grammarlens}

\subsection*{Your turn}
Say these aloud:

\begin{itemize}
\raggedright
  \item \emph{pā̃cvā̃} — two nasal hooks
  \item the rule — \emph{pā̃c} + \emph{-vā̃}
  \item \emph{bīs} $\to$ \emph{bīsvā̃}, twentieth
\end{itemize}

\begin{tcolorbox}[breakable,colback=teal!4,colframe=teal!35!black,title={Before you move on}]
How is an ordinal made from five upward? (\textbf{The plain number, plus -\hi{वाँ}}.) How many nasal hooks are in \emph{\hi{पाँचवाँ}}? (\textbf{Two}.) And say \textquotedblleft{}twentieth\textquotedblright{}. (\textbf{\hi{बीसवाँ}}.)
\end{tcolorbox}
